% =====================================================================
%  Aufgabe 3
% =====================================================================
\chapter{Usability- und UX-Ziele und geplante Anpassungen}
\label{chap:formative_evaluation}

\section{Usability/UX-Ziele}
\label{sec:pt-usability/ux-ziele}

\subsection*{Übergeordnetes Usability- und UX-Ziel (ÜZ)}

\textbf{Studierende sollen ihre individuellen und gemeinsamen Aufgaben in tick.it intuitiv, übersichtlich und mit geringem Aufwand organisieren können.}

\medskip
\noindent
\textit{Zielkategorien: Qualitativ, Ease of Learning und Ease of Use, Performanz und Wahrnehmung.}

\medskip

Das übergeordnete Ziel beschreibt die angestrebte Nutzungsqualität von tick.it und ist daher qualitativ formuliert. Es umfasst sowohl die leichte Erlernbarkeit der Anwendung als auch eine effiziente Bedienung bei regelmäßiger Nutzung. Neben der tatsächlichen Aufgabenbearbeitung soll die Anwendung von den Nutzenden als intuitiv und übersichtlich wahrgenommen werden.

Die Zielsetzung leitet sich aus mehreren Ergebnissen der formativen Evaluation ab. Insbesondere die einfache Bedienung wurde positiv bewertet und soll bei der Weiterentwicklung erhalten bleiben. Gleichzeitig zeigte sich Verbesserungsbedarf bei der Übersicht über anstehende Aufgaben, der Verwaltung von Abgabeterminen und der Organisation von Gruppenarbeiten. Auch der Wunsch nach einer individuellen Strukturierung verschiedener Aufgabenbereiche verdeutlicht die Bedeutung einer übersichtlichen Aufgabenverwaltung (vgl. Abschnitt 2.3).

Aus diesem übergeordneten Ziel werden drei untergeordnete Ziele abgeleitet, welche die Übersichtlichkeit, die intuitive Bedienung und die nachvollziehbare Zusammenarbeit konkretisieren.


\subsection*{Untergeordnetes Ziel 1 (UZ1): Übersichtlichkeit und Priorisierung}

\textbf{Studierende sollen anstehende Aufgaben und Abgabetermine sowie deren Dringlichkeit schnell und eindeutig erkennen können.}

\medskip
\noindent
\textit{Zielkategorien: Quantitativ, Ease of Use, absolut, Performanz.}

\medskip

Das erste untergeordnete Ziel ergibt sich insbesondere aus dem Thema „Deadline-Management und Priorisierung“. Während der Evaluation wurde deutlich, dass die Verwaltung mehrerer gleichzeitig anstehender Aufgaben und Abgabetermine eine Herausforderung darstellen kann. Die befragte Person äußerte den Wunsch nach einer Priorisierungsfunktion und vermisste auch während der Nutzung von tick.it eine Möglichkeit, dringende Aufgaben entsprechend zu kennzeichnen (vgl. Transkript VP01, Z. 122--124).

Darüber hinaus verdeutlicht das Thema „Individuelle Aufgabenorganisation“, dass eine übersichtliche Unterscheidung verschiedener Aufgabenbereiche gewünscht wird. Beispielsweise wurden farbliche Kennzeichnungen für unterschiedliche Studienmodule vorgeschlagen (vgl. Transkript VP01, Z. 160--163).

Aus diesen Erkenntnissen ergibt sich das Ziel, wichtige Aufgaben und zugehörige Abgabetermine ohne aufwendige Suche erkennen zu können. Der Schwerpunkt liegt dabei auf der effizienten Nutzung der Aufgabenübersicht. Das Ziel wird deshalb der Kategorie \textit{Ease of Use} zugeordnet und anhand der benötigten Bearbeitungszeit quantitativ überprüft.


\subsection*{Untergeordnetes Ziel 2 (UZ2): Intuitive Bedienung}

\textbf{Studierende sollen die grundlegenden Funktionen von tick.it bei der erstmaligen Nutzung selbstständig und ohne zusätzliche Anleitung verwenden können.}

\medskip
\noindent
\textit{Zielkategorien: Quantitativ, Ease of Learning, absolut, Performanz.}

\medskip

Das zweite untergeordnete Ziel basiert auf dem Thema „Einfache und intuitive Bedienung“. Im Rahmen der Evaluation wurde insbesondere die unkomplizierte Erstellung neuer Aufgaben positiv hervorgehoben (vgl. Transkript VP01, Z. 143).

Gleichzeitig zeigten sich kleinere Unsicherheiten bei der Bedienung. So war zunächst nicht eindeutig erkennbar, ob das Beschreibungsfeld verpflichtend ausgefüllt werden muss (vgl. Transkript VP01, Z. 145--147). Zusätzlich wurde bei der Datumseingabe eine automatische Vorauswahl des aktuellen Jahres als mögliche Verbesserung genannt (vgl. Transkript VP01, Z. 107--110).

Diese Beobachtungen verdeutlichen, dass die grundlegende Bedienung von tick.it bereits weitgehend verständlich ist, einzelne Gestaltungselemente jedoch noch verbessert werden können. Da im Rahmen der Weiterentwicklung zusätzliche Funktionen vorgesehen sind, soll insbesondere die einfache Erlernbarkeit der Anwendung erhalten bleiben.

Das Ziel konzentriert sich deshalb auf \textit{Ease of Learning}. Es soll überprüft werden, ob Studierende ohne Vorerfahrung grundlegende Aufgaben selbstständig und erfolgreich durchführen können.


\subsection*{Untergeordnetes Ziel 3 (UZ3):
Nachvollziehbare Zusammenarbeit}

\textbf{Studierende sollen bei der gemeinsamen Aufgabenorganisation leicht nachvollziehen können, welche Aufgaben anstehen, wer dafür verantwortlich ist und bis wann diese erledigt werden müssen.}

\medskip
\noindent
\textit{Zielkategorien: Quantitativ, Ease of Learning, absolut, Performanz und Wahrnehmung.}

\medskip

Das dritte untergeordnete Ziel leitet sich aus dem Thema „Kollaborative Aufgabenorganisation“ ab. Die befragte Person beschrieb, dass die Abstimmung innerhalb von Gruppenarbeiten zusätzlichen Kommunikationsaufwand verursachen kann, insbesondere wenn sich Gruppenmitglieder gegenseitig an anstehende Abgabetermine erinnern müssen (vgl. Transkript VP01, Z. 67--73).

Als mögliche Unterstützung wurden gemeinsam nutzbare Boards sowie die Unterteilung größerer Aufgaben in kleinere Arbeitsschritte vorgeschlagen. Diese sollten einzelnen Gruppenmitgliedern zugewiesen und mit individuellen Deadlines versehen werden können (vgl. Transkript VP01, Z. 164--180).

Daraus ergibt sich das Ziel, Verantwortlichkeiten und Abgabetermine bei der gemeinsamen Aufgabenorganisation transparent und verständlich darzustellen. Nutzende sollen die relevanten Informationen auch ohne vorherige Erfahrung mit den Kollaborationsfunktionen erkennen können.

Neben der tatsächlichen Nachvollziehbarkeit soll untersucht werden, wie einfach die gemeinsame Aufgabenorganisation wahrgenommen wird. Das Ziel berücksichtigt deshalb sowohl die \textit{Performanz} als auch die \textit{Wahrnehmung} der Nutzenden.


\section{Geplante Überprüfung der Usability- und UX-Ziele}

Zur Überprüfung der drei untergeordneten Usability- und UX-Ziele sind nach der Umsetzung der entsprechenden Anpassungen aufgabenbasierte Usability-Tests mit Studierenden vorgesehen. Dabei sollen unterschiedliche Aspekte der Nutzung untersucht werden: die Geschwindigkeit beim Erkennen dringender Aufgaben, die erfolgreiche Bedienung ohne Vorerfahrung sowie die Nachvollziehbarkeit der gemeinsamen Aufgabenorganisation.

Für jedes Ziel werden eine geeignete Messmethode, eine Metrik beziehungsweise Skala, ein numerisches Mindestkriterium sowie der Zeitpunkt der Messung festgelegt (siehe Tabelle~\ref{tab:uux_messung}).

Die drei untergeordneten Ziele werden als quantitative und absolute Ziele operationalisiert. Es werden somit feste Mindestkriterien verwendet, ohne die Ergebnisse mit einer anderen Anwendung oder einer früheren Version von tick.it vergleichen zu müssen.


\begin{table}[htbp]
\centering
\caption{Geplante Messung der untergeordneten Usability- und UX-Ziele}
\label{tab:uux_messung}

\footnotesize
\setlength{\tabcolsep}{4pt}
\renewcommand{\arraystretch}{1.3}

\begin{tabularx}{\textwidth}{
    @{}
    >{\raggedright\arraybackslash}p{0.13\textwidth}
    >{\raggedright\arraybackslash}X
    >{\raggedright\arraybackslash}X
    >{\raggedright\arraybackslash}p{0.18\textwidth}
    @{}
}
\toprule
\textbf{Ziel}
&
\textbf{Messmethode und Metrik bzw. Skala}
&
\textbf{Numerisches Mindestkriterium}
&
\textbf{Zeitbezug}
\\
\midrule

\textbf{UZ1: Übersichtlichkeit und Priorisierung}
&
Aufgabenbasierter Usability-Test mit Zeitmessung (\textit{Time on Task}).
&
Median der Bearbeitungszeit maximal 30 Sekunden bis zur korrekten Identifikation der dringendsten Aufgabe und ihres Abgabetermins.
&
Nach einer kurzen Einarbeitung in die überarbeitete Aufgabenübersicht.
\\
\midrule

\textbf{UZ2: Intuitive Bedienung}
&
Aufgabenbasierter Usability-Test; Messung der Aufgabenerfolgsquote ohne Hilfestellung.
&
Mindestens 80\,\% der Testpersonen schließen alle vorgegebenen grundlegenden Aufgaben korrekt und ohne Hilfe ab.
&
Bei der erstmaligen Nutzung der überarbeiteten Anwendung, ohne vorherige Schulung.
\\
\midrule

\textbf{UZ3: Nachvollziehbare Zusammenarbeit}
&
Aufgabenbasierter Usability-Test; Anzahl korrekter Zuordnungen und anschließende \textit{Single Ease Question} (SEQ, 7-Punkte-Skala).
&
Median von mindestens 5 aus 6 korrekten Zuordnungen sowie SEQ-Median von mindestens 5 aus 7 Punkten.
&
Bei erstmaliger Nutzung der Kollaborationsfunktionen; SEQ unmittelbar nach der Testaufgabe.
\\

\bottomrule
\end{tabularx}
\end{table}

\FloatBarrier


Für \textbf{UZ1} wird ein vorbereitetes Kanban-Board mit sechs Aufgaben verwendet, die unterschiedliche Prioritäten und Abgabetermine besitzen. Die Testpersonen sollen die dringendste Aufgabe und deren Abgabetermin identifizieren.

Erfasst wird die Zeit vom Anzeigen des Boards bis zur korrekten Antwort. Eine Aufgabe gilt nur bei korrekter Identifikation als erfolgreich abgeschlossen. Erfolgt innerhalb von 90 Sekunden keine korrekte Antwort, wird der Durchgang als nicht erfolgreich gewertet und für die Berechnung des Medians mit 90 Sekunden berücksichtigt.

Der angestrebte Median von maximal 30 Sekunden dient als vorläufiger Richtwert für eine schnelle Orientierung in einer überschaubaren Aufgabenliste. Da das Ziel die effiziente Bedienung betrifft, erfolgt die Messung nach einer kurzen Einarbeitung in die grundlegende Nutzung der Aufgabenübersicht.

Für \textbf{UZ2} bearbeiten die Testpersonen bei ihrer ersten Begegnung mit der überarbeiteten Anwendung drei grundlegende Aufgaben: das Erstellen einer Aufgabe mit Abgabetermin, das nachträgliche Bearbeiten einer bestehenden Aufgabe und das Verschieben einer Aufgabe zwischen den Kanban-Spalten.

Erfasst wird, wie viele Testpersonen sämtliche Aufgaben korrekt und ohne Hilfestellung abschließen können. Die angestrebte Erfolgsquote von mindestens 80\,\% soll sicherstellen, dass die grundlegende Bedienung für einen großen Teil der Testpersonen ohne vorherige Schulung verständlich ist.

Da die leichte Erlernbarkeit untersucht wird, erhalten die Testpersonen vor Beginn keine Einführung in die konkreten Bedienabläufe.

Für \textbf{UZ3} wird ein gemeinsames Board mit drei vorbereiteten Aufgaben verwendet. Die Testpersonen sollen für jede Aufgabe die zuständige Person und den zugehörigen Abgabetermin identifizieren. Daraus ergeben sich insgesamt sechs überprüfbare Zuordnungen.

Als Mindestkriterium wird festgelegt, dass im Median mindestens fünf dieser sechs Informationen korrekt erkannt werden. Zusätzlich beantworten die Testpersonen unmittelbar nach der Aufgabe eine \textit{Single Ease Question} (SEQ), mit der die wahrgenommene Einfachheit der Aufgabenbearbeitung auf einer siebenstufigen Skala von sehr schwierig bis sehr einfach bewertet wird.

Dabei soll ein Median von mindestens fünf Punkten erreicht werden. Die Kombination aus korrekten Zuordnungen und subjektiver Bewertung ermöglicht es, sowohl die tatsächliche Verständlichkeit als auch die wahrgenommene Einfachheit der gemeinsamen Aufgabenorganisation zu untersuchen.

Da die bisherige formative Evaluation auf einem Interview basiert, sind die abgeleiteten Ziele zunächst als begründete Gestaltungshypothesen zu betrachten. Die geplanten Usability-Tests sollen daher überprüfen, inwieweit die Weiterentwicklung von tick.it die identifizierten Bedürfnisse unterstützt und ob die festgelegten Ziele erreicht werden.

% ---------------------------------------------------------------------
\begin{beitraege}
    Interview Vorbereitung & Tara Dethlefsen \\
    Interview Durchführung & Philipp Wisser \\
    Interview Transkription & Philipp Wisser \\
    Thematische Analyse in MAXQDA & Lara Senger \\
  Text \cref{sec:pt-einleitung2} und  \cref{sec:pt-vorgehen} & Tara Dethlefsen \\
  Text \cref{sec:pt-ergebnisse}   & Lara Senger \\
  \midrule
  KI/Automatisierung: Claude Sonnet 5.5 zur Überprüfung von Grammatik und Rechtschreibung \parencite{claude2026sonnet}  & Tara Dethlefsen \\
  KI/Automatisierung: Whisper (Modell medium) zur automatischen Transkription des Interviews \parencite{whisper2025}  & Philipp Wisser \\
  KI/Automatisierung: Claude Opus 5.5 zur Prüfung der Formvorgaben \parencite{claude2026review} & Philipp Wisser \\
\end{beitraege}
